\section{Bounded search among tied minimum-span surfaces}
\label{sec:cocycle-face}

The minimum number of normal pieces does not determine the topology of a
cocycle surface. Even a fixed optimal dual matching can certify both a disc
and a non-disc on the same source exterior. This distinction motivates a
new optional bounded fallback: after the selected optimum and all requested tree
gauges miss, inspect several other potentials with the same minimum span.
The construction uses shortest paths in a difference-constraint graph;
it does not run another flow optimization. Every positive retains the
existing independent source-bound certificate format.

\subsection{An exact description of all tied optima}

Fix a verified span certificate with optimum $D_*$ and matching rows
$(t,i,s,j)$. Let $\ell_t$ be its selected lower corner in tetrahedron $t$,
and let $u_t$ be its selected upper corner. Each exists exactly once because
the matching uses every lower and upper tetrahedron once. For any potential
$f$, write $a_{t,i}=h_{t,i}+f(v_{t,i})$. Cancellation at matched global
vertices gives
\[
 D(f)-D_* = \sum_t \bigl(\max_i a_{t,i}-a_{t,u_t}\bigr)
              +\sum_t\bigl(a_{t,\ell_t}-\min_i a_{t,i}\bigr).
\]
Every summand is nonnegative. Consequently $f$ is optimal if and only if
all selected lower and upper corners remain extremal. Equivalently, for
each tetrahedron and each of its four corners,
\begin{align*}
 f(v_{t,\ell_t})-f(v_{t,i})&\leq h_{t,i}-h_{t,\ell_t},\\
 f(v_{t,i})-f(v_{t,u_t})&\leq h_{t,u_t}-h_{t,i}.
\end{align*}
These eight inequalities per tetrahedron describe \emph{every} optimum,
including those admitting other optimal matchings. Necessity follows from
weak-duality equality with this one fixed optimal matching; no uniqueness
assumption is involved.

Represent $f(b)-f(a)\leq w$ by an arc $a\to b$ of weight $w$.
Given the certificate's original potential $f_0$, reweight it as
$r(a,b)=w+f_0(a)-f_0(b)\geq0$. In displacement variables $q=f-f_0$,
the inequalities become $q(b)-q(a)\leq r(a,b)$. The graph has at most
$4T$ vertices and exactly $8T$ arcs, retaining parallel arcs and loops.
The weights and the constructed extremal solutions are integral.

Within each connected component of the tetrahedron--vertex incidence
structure, this constraint graph is strongly connected. Indeed, fixing
$D(f)=D_*$ bounds every local height span by $D_*$. Any two corners of a
tetrahedron therefore have potential difference at most
$D_*+|h_{t,i}-h_{t,j}|$. Chains of tetrahedra bound every relative potential
in that incidence component. If the directed constraint graph were not
strongly connected, decreasing all potentials in a sink strongly connected
component would preserve all inequalities indefinitely, contradicting that
bound. Different incidence components can translate independently.

\subsection{Two simultaneous extrema from one root}

For a root $r$ in a component, let $d(r,v)$ be shortest-path distance in
the nonnegative reweighted graph. Normalize the displacement by $q(r)=0$.
Every feasible displacement obeys
\[
 -d(v,r)\leq q(v)\leq d(r,v).
\]
Both simultaneous bound vectors are feasible:
\[
 q^+(v)=d(r,v),\qquad q^-(v)=-d(v,r).
\]
The upper assertion is the shortest-path triangle inequality along an arc;
the lower assertion follows from the same inequality for paths ending at
$r$. Thus each vector attains all its corresponding coordinate bounds at
once. One Dijkstra search in each direction constructs the pair. Outside
the root's incidence component, the producer keeps $q=0$. A common additive
normalization then makes the smallest output potential zero.

If two roots have zero distance in both directions, they are in the same
zero-weight strongly connected component. Their forward distance vectors
are identical, as are their reverse distance vectors. Testing both roots
would therefore produce the same pair. The implementation marks all such
roots after the two searches and skips them. It orders the remaining roots
by descending total arc incidence, breaking ties by vertex label. It also
skips normalized potentials already emitted or equal to the original one.
These are exact duplicate rules, not assumptions about surface topology.

Coordinates are reconstructed from the adjusted tetrahedral heights, and
every emitted potential replays the unchanged span verifier with the
original matching and optimum. At the source stage, the original primitive
class and minimum-span connectedness theorem still apply. Negative
candidates use normal matching and Euler accounting; a proposed disc must
also pass the complete independent source-bound verifier. Failure to find
a disc remains inconclusive.

\subsection{Cost and scheduling limits}

Let $R$ be the smaller of the requested root allowance and the number of
vertices. After verification and graph construction, at most $2R$
shortest-path searches and $2R$ candidate replays occur. In an indexed-graph
model with constant-time table lookup and well-distributed signature hashes,
the arithmetic cost is $O((R+1)T\log(T+1))$, including label/root sorting.
This is a conditional container-cost bound, not an unconditional worst-case
bound for Python's hash tables.

The shipped code uses dictionaries for labels and distances and a set of
potential tuples for duplicates. Adversarial hash collisions can require
$O(T)$ probes per dictionary operation and $O(RT)$ coordinate comparisons
per signature lookup. Since $R\leq4T$, a conservative deterministic bound
for these container comparisons together with the graph arithmetic is
$O((R+1)T^2)$. Storage is $O((R+1)T)$ integer records; stored duplicate
signatures account for the factor $R$. With $B$ the largest input integer
bit length, intermediate quantities use $O(B+\log(T+1))$ bits. Charging
linear bit cost for integer hashing, addition and comparison gives the
corresponding product bound. No loop ranges over a height's numerical
value. Both bounds are local polynomial bounds, including when all roots
are requested.

The optional source-stage experiment requests four roots, after all
previously requested trees and the single span optimization have failed.
This placement preserves earlier positive proofs and their guard work.
The shared cap and cancellation callback cover the new graph searches,
coordinate checks and final replay. Setting \code{face\_roots=0}, or
\code{normal\_seed\_face\_roots=0} in the recognizer, restores the
previous search policy. The CLI exposes \code{--normal-seed-face-roots}.
Without span optimization the face fallback is skipped. The face allowance defaults to zero: the measurements below show added
coverage but slower complete calls. The entire native normal-surface stage
also remains optional in the ordinary portfolio.

Even trying every root does not enumerate all feasible potentials, all
vertices of the optimality polytope, or all topology types. The exact
inequality description is complete; this extremal candidate search is not.
It does not maximize Euler characteristic, prove minimal genus, or complete
the normal-surface search required for general quasipolynomial recognition.

\subsection{Topology changes at an unchanged optimum}

On the five-crossing source \code{optimized-positive}, the original
84-piece optimum is a disc. Another potential certified by the very same
matching has 84 pieces and is an annulus. Independent source inspection
and Regina agree that both surfaces are connected and orientable; their
Euler characteristics are 1 and 0, and their boundary counts are 1 and 2.
A dual matching therefore does not fix the normal coordinates or topology.
A change in flow tie-breaking must not be justified by assuming that it
preserves the selected surface.

The seven-crossing source \code{random-142} supplies the converse discovery
example. Its original optimum has 122 pieces and is an annulus. Four
nonequivalent root trials find a 122-piece disc at root 13 in the forward
direction, after five distinct non-disc alternatives. An exhaustive pass
through root classes also finds a second disc and two three-boundary planar
surfaces, all at the same optimum. The retained witnesses share the original
matching; they differ in potential and coordinates. The successful bounded
source proof is independently replayed with the discovery entry points
disabled in a focused test.

With four face roots requested, native positives on the 84-source corpus
increase from 23 to 24 under the four-tree policy, or from 24 to 25 when
24 tree trials are requested. There is one distinct new positive proof.
The shared-budget Gordian cap remains, leaving 59 and 58 completed restricted
misses respectively. The new seven-crossing proof uses 117,848 guards under
the four-tree policy; disabling face search leaves a completed miss at
93,011 guards. Every previously positive certificate and its guard work are
preserved. These counts describe the explicit four-root option; the default
zero-root policy keeps its previous coverage.

\subsection{Validation and the decision to keep this optional}

The 32 focused tests pass in 3.991 seconds. They include forty small
arithmetic instances checked on exhaustive potential grids, malformed
optimality witnesses, disconnected incidence components, 20,000-bit scaled
heights, cancellation, strict option types and exact positive-proof budget
boundaries. The final maintained suite passes all 1,197 tests in
232.354 seconds. An earlier full run before changing the default allowance
also passed; the final run covers the disabled-by-default policy.

The source audit independently inspects every distinct root-extremal
surface on the 74 inputs with at most twelve crossings, together with each
original optimum. All 787 surfaces agree with Regina on connectedness,
orientability, Euler characteristic and compressing-disc status. There are
713 alternatives, and 21 of the 74 sources exhibit different Euler
characteristics among their tied optima. Full witnesses for the two named
counterexamples and the new positive are retained. On nineteen ordinary
portfolio inputs, both zero-root and four-root settings return the expected
verdicts without reaching the native stage.

The first audit reached the final ordinary-corpus loop but failed because
its driver used an incorrect corpus-helper attribute. The helper reference
was corrected and the failure log retained. The final pinned audit takes
103.135 seconds, overlapping correctness tests; this is not a performance
comparison. Its 485 runtime, driver, fixture and prior-audit hashes agree
with the final benchmark. The independent span and source-verifier modules
are byte-identical to the previous baseline \code{a8658ec40}.

Two serial five-round experiments compare the complete recognizer with
face search explicitly off or set to four roots. Each case has one warm-up
round and then five measured rounds, each randomly ordering off, off A/A,
on and on A/A calls. No tests, audits, builds or commits overlap the timed
runs. Earlier shortcuts are disabled so that the native stage is reached;
both arms have no local seed cap and retain the twelve-second global
recognition deadline. All requested tree trials remain the same in both
arms, including 24 for the late-tree control. A returned native proof is
additionally replayed and hashed inside the timed endpoint.

The first experiment motivated the zero default. The second repeats the
comparison with the final source pins. All 280 measured calls and 56
warm-ups across the two experiments finish; per-implementation verdicts
and proof hashes are stable across both runs. The final run alone has
140 measured calls and 28 warm-ups and takes 15.003 seconds. Its primary-arm
median times and median paired ratios are shown below; a median paired
ratio need not equal the ratio of median times.

\begin{center}
\small\begin{tabular}{lrrrrr}
\toprule
Input & Off (ms) & On (ms) & Off/on & Off A/A & On A/A \\
\midrule
Raw circle & 63.764 & 64.729 & 0.986 & 1.003 & 1.011 \\
Optimized positive & 92.861 & 91.446 & 1.015 & 1.002 & 1.003 \\
Early tree positive & 37.904 & 38.139 & 0.992 & 0.951 & 1.020 \\
Late tree positive & 140.885 & 136.867 & 1.028 & 1.065 & 1.007 \\
New face positive & 120.063 & 206.335 & 0.582 & 0.955 & 0.973 \\
Trefoil & 41.078 & 64.424 & 0.638 & 0.984 & 1.025 \\
Genus-one miss & 43.752 & 66.204 & 0.667 & 0.993 & 1.000 \\
\bottomrule
\end{tabular}

\end{center}

The new seven-crossing native positive takes about 206 milliseconds with
face search versus 120 milliseconds when the existing exact fallback
finishes the call. Its paired ratio corresponds to about 72 percent
additional time. The two continuing misses cost about 50 and 57 percent
more. Earlier positive branches avoid the new search; their small timing
differences provide no evidence of algorithmic gains and are accompanied
by visible A/A variation. In particular, the late-tree off-arm A/A ratio
is 1.065. The experiment therefore supports optional discovery and stronger
coverage evidence, not a faster default recognizer. The maintained default
continues to request zero face roots.

Reproduction uses \path{../fast/normal_orbit_research/face.py} with
\code{audit} or \code{recognize}; the audit requires optional Regina.
The existing incoming archive inventory is unchanged. The checked
\code{openai/math} head remains \code{fd4aeeb2e}, and the arXiv record for
Lackenby's hierarchy paper~\cite{lackenby2026} still lists its July 2026 first version. The
new optimal-face argument is a local extension of report 31's matching
duality, not a resolution of the general recognition-complexity gap.


